\section{A rational cubic seven-factor on the same mixed window}
\label{sec:mixed-cubic-seven-window}

\paragraph{The preceding original-coefficient remainder.}
Retain every datum, page, prescribed pair, field, original coefficient,
lattice, normalization and full local condition of AC-G208-window.
In particular $h\geq2$ and $T\geq2$ are integers, $N=T-1$, and
\[
 \begin{gathered}
 H_u=\mathbb Q(\zeta_{2^{h+2}},\zeta_7,\sqrt{-q},\sqrt{-d},
                  \omega,\sqrt[3]{269}),\qquad K=\mathbb Q(\omega),\\
 u^{2+8\varepsilon}\leq2^h\leq2u^{2+8\varepsilon},\qquad
 T=\lfloor u^{399}\rfloor,\qquad0<\varepsilon<1/256.
 \end{gathered}\tag{MC7W.1}
\]
Keep exactly $Z_u,J_{u,0},J_{u,\rm mix},V_u,Y_u$,
$C_5^{\rm all}(u),U_5^{\rm all}(u),C_7^{\rm all}(u),U_7^{\rm all}(u)$
and $M_8(u)$ of (MQW.2)--(MQW.6).
Write $\chi$ for either faithful Kummer cubic character and
$\nu=\xi_7^2$ for either faithful rational cubic seven-character.
The two cubic factors are independent in (FC.1).

For primes $\ell>3$, $\ell\ne7,269$, use precisely the
indicators of AC-G209 (MC7.2): $c(\ell)=1$ if and only if
$\ell\equiv1\pmod3$ and $269$ is a noncube modulo $\ell$;
otherwise $c(\ell)=0$. Let $s(\ell)=0$ if
$\ell\equiv\pm1\pmod7$ and $s(\ell)=1$ otherwise.
Define
\[
 \begin{split}
 \mathscr P_7&=\{\ell\text{ prime}:\ell>3,\ \ell\ne7,269,
                                      \ c(\ell)+s(\ell)=1\},\\
 M_7(u)=\{(\ell,\chi^j\xi_7^{2a},e):{}&\ell\in\mathscr P_7,
                         \quad j,a\in\{1,2\},\\
                         &1\leq e<T,\quad(\ell-1)\mid e\},\\
 J_7(T)&=\sum_{\substack{\ell\in\mathscr P_7\\\ell\leq T}}
                       \left\lfloor\frac{N}{\ell-1}\right\rfloor,
 \qquad M_{8,7}(u)=M_8(u)\cup M_7(u).
 \end{split}\tag{MC7W.2}
\]
Every undisplayed two-power and auxiliary $q,d$ factor of a
new label is trivial. The exclusive condition in $\mathscr P_7$
is the full stated criterion; it is not restricted to split primes
or to the older set $\mathscr P_{\rm cub}$.

Define
\[
 \begin{split}
 \mathcal R_{399}^{\rm match,cub8,7,5all}(u)={}&
  \sum_{(\ell,\eta,e)\in J_{u,\rm mix}\setminus M_{8,7}(u)}
                      \frac{\log\#\mathcal S_{\ell,\eta}(e)}{b_\ell}\\
 &+\sum_{(\ell,\eta,e)\in
      J_{u,0}\setminus(V_u\cup Y_u\cup C_7^{\rm all}(u)
                                      \cup C_5^{\rm all}(u))}
       \frac{\log\#\mathcal S^{\mathbb Q}_{\ell,
                           \psi_{\eta,e}^{\rm nc}}(e)}{b_\ell}.
 \end{split}\tag{MC7W.3}
\]
These are the same finite nonnegative sums as in (MQW.3),
with the indicated additional full mixed terms removed.

\begin{theorem}[AC-G210-window: the full exclusive local criterion]
\label{AC-G210-window}
Every triple of $M_7(u)$ has zero full original $K$-strict group.
It is a new retained mixed triple, with the following complete
old-overlap and pointwise remainder identities:
\[
 \begin{gathered}
 M_7(u)\subset J_{u,\rm mix}\setminus M_8(u),\\
 M_7(u)\cap\bigl(Z_u\cup V_u\cup Y_u\cup C_7^{\rm all}(u)
                            \cup C_5^{\rm all}(u)\cup M_8(u)\bigr)
                                            =\varnothing,\\
 M_{8,7}(u)=M_8(u)\sqcup M_7(u),\\
 \mathcal R_{399}^{\rm match,cub8,7,5all}(u)
                     =\mathcal R_{399}^{\rm match,cub8,5all}(u).
 \end{gathered}\tag{MC7W.4}
\]
The exact new count and an explicit lower bound are
\[
 \begin{gathered}
 |M_7(u)|=4J_7(T),\qquad
 |M_{8,7}(u)|=8J_{\rm cub}(T)+4J_7(T),\\
 |M_7(u)|\geq4\lfloor N/4\rfloor+4\lfloor N/10\rfloor
                                      +4\lfloor N/12\rfloor.
 \end{gathered}\tag{MC7W.5}
\]
This is a proper reduction of the retained indexing set for $T\geq5$.
The full disjoint represented-triple omission and its cardinality are
\[
 \begin{gathered}
 Z_u\sqcup V_u\sqcup Y_u\sqcup U_7^{\rm all}(u)
                 \sqcup U_5^{\rm all}(u)\sqcup M_{8,7}(u),\\
 |Z_u|+|V_u|+|Y_u|+47N-\lfloor N/4\rfloor
                            +8J_{\rm cub}(T)+4J_7(T).
 \end{gathered}\tag{MC7W.6}
\]
With exactly the preceding constants and costs,
\[
 \begin{aligned}
 |I_{399}-\mathcal R_{399}^{\rm match,cub8,7,5all}|&\leq
        \mathcal E_u+\mathcal L_u+\mathcal A_u+C_2\mathcal P_u,\\
 |I_{399}^{\rm all7,2}-\mathcal R_{399}^{\rm match,cub8,7,5all}|&\leq
        \mathcal E_u+\mathcal L_u+C_2\mathcal P_u.
 \end{aligned}\tag{MC7W.7}
\]
There is no additional local, torsion, critical or family-count cost.
The same envelope $O_\varepsilon(u^{806+16\varepsilon}(\log u)^2)$,
the identical common window, and the preceding threshold/slack
implications remain valid.
\end{theorem}

\begin{proof}
\proofstep{1. Actual mixed labels and the exact original coefficient.}
In (FC.1) the Kummer $C_3$ and the cubic subfactor of the
rational conductor-seven $C_6$ are independent direct factors.
Their two nontrivial powers give four distinct products
$\chi^j\xi_7^{2a}$. All these characters already occur in
the host (MC7W.1). Their values lie in the prescribed original
completion $O_\ell$ of $\mathbb Z[\zeta_{3\cdot2^h}]$,
since only $\mu_3$ is required.
Retain its fraction field $E_\ell$, embedding, uniformizer
and normalizing rank $b_\ell$. The coefficient of (FC.2) is
\[
 (O_\ell\otimes_{\mathbb Z_\ell}\mathbb Q_\ell/\mathbb Z_\ell)
                       (\chi^j\xi_7^{2a})(-e)
       =(E_\ell/O_\ell)(\chi^j\xi_7^{2a}c_\ell^{-e}).
 \tag{MC7W.8}
\]
Apply AC-G209 (MC7.3)--(MC7.4) with $\kappa=1$.
The full prime and weight conditions are exactly (MC7W.2).
G209 Step5 supplies the original-host inflation from
$G_{K,\Sigma_\ell}$, preserving every full completion and
away-inertia condition. Thus the original group
$\mathcal S_{\ell,\chi^j\xi_7^{2a}}(e)$ is zero.

This application includes both local alternatives. If $c=1,s=0$,
the prime splits in $K$, the Kummer Frobenius has order three
at both completions, and the rational cubic factor is locally trivial.
If $c=0,s=1$, the Kummer cubic splits locally, while the rational
cubic value has order three after the residue-degree power
$[K_v:\mathbb Q_\ell]=1$ or two. This includes the unique
inert completion as well as both completions of a split prime.
The exclusive criterion makes the product nontrivial at every
coefficient completion for all four character pairs.
No second split local condition is omitted and no split-prime
restriction is imposed on the complementary case.

The class input is precisely the degree-nine quotient in
(MC7.11), $\operatorname{Cl}(B)$ a quotient of $C_{14}\times C_2$.
The transposition argument kills the required nontrivial Kummer
components in the degree-eighteen cover, irrespective of the
rational cubic component; it does not prove the entire larger
class group zero. A displayed character has order three,
not nine. The degree-nine cover over $K$ trivializes all four
characters and need not be their minimal splitting field.
These distinctions and the original-uniformizer lift are the
literal G209 input, with no new class or unit computation here.
The finite-etale coefficient and strict-kernel dictionary on
$\operatorname{Spec}\mathcal O_K[1/(3\cdot7\cdot269\cdot\ell)]$
is retained from G209; no new archival or motivic claim is made.

\proofstep{2. Exhaustive old intersections in both prime alternatives.}
Every new label has nontrivial Kummer factor. The old $Z_u$
mixed omissions are confined to coefficient $269$, excluded
in (MC7W.2). The fixed coefficient-two and coefficient-three
families are excluded as well. Every other old $Z_u$ family,
including the growing high-two and inert-trivial rows, has
trivial Kummer factor and cannot contain a new label.
This remains true at coefficient five, even when the old
inert-trivial row has the same weight divisible by four.

The sets $V_u,Y_u,C_7^{\rm all}(u),C_5^{\rm all}(u)$ are
rational families with trivial Kummer factor. In particular,
the rational cubic seven-powers already present at coefficient
five are different from the new products with $\chi^j$.
The same independence distinguishes every auxiliary-prime row
in $Y_u$ if a new coefficient happens to equal $q$ or $d$.
Thus every new row lies in $J_{u,\rm mix}$.

All eight labels of $M_8(u)$ have trivial rational seven-factor.
All four labels of $M_7(u)$ have nontrivial rational cubic
seven-factor. Their intersection is therefore empty, even
when their eligible prime and weight coordinates agree.
The overlap of the two prime sets does not duplicate triples.
This proves (MC7W.4) and the disjoint union (MC7W.6), keeping
every previous mixed or noncritical family unchanged.

\proofstep{3. Exact broad-set count and pointwise deletion.}
For each prime in $\mathscr P_7$, the number of positive
weights below the integer $T$ divisible by $\ell-1$ is
$\lfloor N/(\ell-1)\rfloor$. If nonzero, it implies $\ell\leq T$.
This is derived from the weight constraint and does not replace
the unrestricted coefficient-prime domain of (TC.1).
The four distinct characters give $4J_7(T)$ new triples.
Together with $|M_8|=8J_{\rm cub}(T)$ this proves (MC7W.5)
and the cardinality in (MC7W.6).

Primes five and eleven are inert in $K$ and not congruent to
$\pm1$ modulo seven, so they have $c=0,s=1$.
At thirteen, $c=1$ because $269^4\equiv9\ne1\pmod{13}$,
while $13\equiv-1\pmod7$ gives $s=0$.
Their separate progressions give the lower bound in (MC7W.5).
For $T\geq5$, $(5,\chi\xi_7^2,4)$ is newly absent.
No density or quantitative prime-distribution premise is used.

Each new summand is the full mixed group proved zero in Step1,
so its normalized value is $\log1/b_\ell=0$.
Delete exactly those terms from (MQW.3), or apply (MC7.7)
to the currently retained triples. This yields (MC7W.3)
and its pointwise equality with the preceding remainder.
The whole rational noncritical sum and every existing critical
comparison remain unchanged. No new critical partner appears,
and no positive derived index is split or declared zero.

\proofstep{4. One inherited cost and the original common window.}
Substitute (MC7W.4) into (MQW.7).
One full (FL.4) local majorant, one (LT.2) torsion aggregate,
one existing critical aggregate and the old G189 full-index
cost give exactly (MC7W.7). There is no multiplier for four,
the prime or weight counts, or the degree of the proof cover.
The original degrees $D_u=144\cdot2^h$ and $d_u=96\cdot2^h$
retain their distinct source roles.

Fix the same datum, page, $0<\varepsilon<1/256$,
parameters $J,X_0,s$, prescribed pairs and actual window
$X^{1/J}\leq W$, $W^2\leq X$ of (TC.3).
Equal remainders have the same numerical premise, threshold
and slack on this identical $W$. For the original full index
retain the same uniform $U_0$, threshold $\max(X_0,U_0^J)$
and half-slack argument. No independent existential windows
are intersected. Every simultaneous original coefficient,
lattice, field, finite condition, actual rank $R.t$, page
and pair remains under the same quantifiers.
\end{proof}

\paragraph{Scope and chronology.}
This is one composition for the full exclusive criterion of G209,
including both its split-noncube branch and its complementary
local branch. No narrower intermediate removal is credited.
Every other mixed and previously retained noncritical term remains
in (MC7W.3). Their numerical bound, mixed theta/Selmer realization
and H04.4 stay open. The G209 arbitrary-set identity and this
explicit old-overlap application are assessed together at their
actual completion times, with no packaging credit or later banking.
No new question row or automatic gate closure is asserted.
